\uuid{F1bX}
\titre{Famille libre/Génératrice/Base}
\auteur{hanani ; mzari}
\organisation{exo7}
\licence{CC BY-NC-SA 4.0}
\datecreate{2020-01-20}
\niveau{L1}
\module{Algèbre}
\chapitre{Espace vectoriel}
\sousChapitre{Base}
\difficulte{3}
\theme{Base et dimension, Famille libre, Génératrice, Base}
\exo7qcmid{lille-2.espaces-vectoriels.039}

\contenu{
\begin{question}
\qclassification{106.04}
\qtags{motcle=Famille libre/Génératrice/Base}
Soit $E$ un espace vectoriel sur $\Rr$ de dimension $3$ et$v_1,v_2,v_3$ des vecteurs linéairement indépendants de $E$.
Quelles sont les assertions vraies ?
\begin{answers}
\good{$\{v_1,v_2,v_3\}$ est une famille génératrice de $E$}.
\good{$\{v_1,v_2,v_1+v_3\}$ est une base de $E$}.
\bad{$\{v_1-v_2,v_1+v_3\}$ est une base de $E$}.
\good{$\{v_1-v_2,v_1+v_3\}$ est famille libre de $E$}.
\end{answers}
\begin{explanations} Puisque $\{v_1,v_2,v_3\}$  est une famille libre qui contient $3$ vecteurs et la dimension de $E$ est $3$,
elle est génératrice et donc c'est une base de $E$.
\vskip2mm
On vérifie aussi que $\{v_1,v_2,v_1+v_3\}$ est une famille libre,
et donc pour les mêmes raisons que précédemment, c'est une base de $E$.
\end{explanations}
\end{question}
}
